\documentclass[10pt]{article}
\providecommand{\FIGDIR}{fig}
\newcommand{\DIGESTDATE}{21 October 2026}
\input{preamble}

\begin{document}

% ==================================================================== MASTHEAD
\thispagestyle{empty}
\begin{tikzpicture}[remember picture, overlay]
  \fill[Deep] (current page.north west) rectangle
      ($(current page.north east)+(0,-67mm)$);
  \fill[Amber] ($(current page.north west)+(0,-67mm)$) rectangle
      ($(current page.north east)+(0,-68.4mm)$);
  \foreach \x/\y/\r in {22/12/0.9, 41/26/1.5, 63/9/0.7, 88/31/1.1, 112/17/1.9,
                        138/29/0.8, 159/13/1.3, 178/54/1.0, 196/21/1.6,
                        30/44/0.6, 74/57/1.2, 125/48/0.9, 168/6/0.8, 191/59/0.7,
                        52/18/1.0, 100/60/0.6, 148/20/1.1, 66/49/1.4, 205/45/0.8}
    \fill[Sky, opacity=0.30] ($(current page.north west)+(\x mm,-\y mm)$) circle (\r mm);
\end{tikzpicture}

\vspace*{4mm}
{\sffamily\fontsize{7.6}{9}\selectfont\color{Sky}%
 AUTOMATED DAILY LITERATURE DIGEST\;\textbullet\;PhD RESEARCH WATCH\par}
\vspace{2.2mm}
{\sffamily\bfseries\fontsize{27}{29}\selectfont\color{white}%
 Particulate Matter\par}
{\sffamily\fontsize{27}{29}\selectfont\color{Sky}%
 Monitoring \& Health Effects\par}

\vspace{5.0mm}
{\sffamily\fontsize{8.4}{10}\selectfont\color{white}%
 Issue 2026-10-21\;\;\textcolor{Amber}{\textbar}\;\;Entry window 21 October 2026 (full day)\;\;%
 \textcolor{Amber}{\textbar}\;\;18 records screened in scope\par}
\vspace{18mm}

\noindent
\begin{minipage}[t]{0.190\linewidth}\metric{18}{RECORDS IN SCOPE\\(86 EXCLUDED)}{Deep}\end{minipage}\hfill
\begin{minipage}[t]{0.190\linewidth}\metric{9}{OF 10 SUBTOPIC\\BANDS POPULATED}{Teal}\end{minipage}\hfill
\begin{minipage}[t]{0.190\linewidth}\metric{5}{EFFECT ESTIMATES\\WITH A CI}{Sky}\end{minipage}\hfill
\begin{minipage}[t]{0.190\linewidth}\metric{3}{TIER-A\\RECORDS}{Amber}\end{minipage}\hfill
\begin{minipage}[t]{0.190\linewidth}\metric{5}{RECORDS CARRIED ON\\METADATA ALONE}{Clay}\end{minipage}

\vspace{6mm}

\section*{\sffamily\bfseries\color{Deep}Signal of the day}
\vspace{-1.5mm}{\color{Amber}\rule{\linewidth}{1.1pt}}\vspace{2.5mm}

\begin{keybox}
{\sffamily\bfseries\small\color{Deep}26 million Medicare beneficiaries: atrial fibrillation tracks sulfate and
combustion PM$_{2.5}$, not biomass smoke or secondary organics.}\\[1.4mm]
\footnotesize
Source-apportioned PM$_{2.5}$ (Multilinear Engine on speciation-network data, inverse-distance weighted to ZIP code)
against 4.44 million incident AF cases over 171.6 million person-years, 2000--2019. Per IQR: secondary sulfate HR
\textbf{1.056} (1.053--1.058), heavy oil/coal combustion \textbf{1.048}, secondary nitrate 1.032, vehicle 1.029;
biomass burning and dust modest; SOA null nationally but positive in the West.
\textbf{Where it is weak:} source fields are interpolated from sparse speciation monitors over 100\,km, so source
contrasts are largely regional contrasts, and per-IQR scaling differs by source --- this is not a per-microgram
toxicity ranking. Claims-based AF misses paroxysmal disease.
\textbf{Why it matters here:} optical low-cost networks report mass. If composition carries the cardiac signal,
a dense mass network needs a sparse speciation backbone to be epidemiologically useful.\par
\end{keybox}

\vspace{3mm}

\section*{\sffamily\bfseries\color{Deep}What changed today}
\vspace{-1.5mm}{\color{Amber}\rule{\linewidth}{1.1pt}}\vspace{2.5mm}

\footnotesize
\begin{itemize}

\item \textbf{Particle size matters for spirometry.} China Pulmonary Health Study, 50,991 adults: per IQR PM$_1$
FEV$_1$ \textbf{$-$36.8\,mL} ($-$45.9 to $-$27.8); large-airway decrements grow as size falls, while the coarse
fraction PM$_{2.5\text{--}10}$ gives the largest small-airway (FEF$_{25\text{--}75}$) decrements.

\item \textbf{An outcome-weighted air index.} 1.21 million AMI cases, 325 cities: a composite Air Health Index
(PM$_{2.5}$, NO$_2$, O$_3$, temperature) gives +5.97\,\% per IQR vs +3.33\,\% for AQHI and +2.04\,\% for AQI ---
in-sample advantage expected, but externally validated.

\item \textbf{Implausibly large cancer-prognosis HRs.} Cervical cancer, 3,758 women: PM$_{2.5}$ recurrence SHR
1.93, death HR 2.55 per IQR on 113 deaths --- read as confounding by referral geography until replicated.

\item \textbf{Magnetite inventory.} Respirable magnetite from Chinese coal power: 2,829\,t\,y$^{-1}$; two provinces
with 4\,\% of output could cut $>$31\,\% of emissions by tightening standards.

\item \textbf{PubMed was nearly silent.} Only 2 records carried a 21 October entry date in PubMed; Europe PMC and
the Crossref ISSN sweep carried the issue, and five ACS \emph{Environment \& Health} papers arrived via Europe PMC
with PMIDs assigned the next day.

\end{itemize}

\vspace{2mm}

\begin{keybox}
{\sffamily\bfseries\footnotesize\color{Deep}Window continuity}\\[1.2mm]
{\fontsize{7.2}{8.8}\selectfont
The 20 October issue closed with \texttt{last\_entry\_date} at 20 October. This issue collects
\textbf{21 October in full}, continuous with it. The scheduled run for 21 October (04:0x UTC 22 October) did not
fire; this run started 18:02 CDT on 22 October, before the 22:00 cut-off, so 22 October is \emph{not} opened and
is left to the next run at or after 22:00 CDT. The W43 weekly (18--24 October) is due Saturday 24 October.\par}
\end{keybox}

\vspace{3mm}
\par\vskip 0pt plus 18\baselineskip\penalty-200\vskip 0pt plus -18\baselineskip
\section*{\sffamily\bfseries\color{Deep}Corpus analytics}
\vspace{-1.5mm}{\color{Amber}\rule{\linewidth}{1.1pt}}\vspace{2.5mm}

\noindent\begin{minipage}{\linewidth}\includegraphics[width=\linewidth]{\FIGDIR/f1_subtopics.png}
\figcap{Subtopic assignment of the 18 in-scope records across nine bands. Cardiovascular, other
clinical endpoints, exposure assessment and burden/mitigation hold three each; sensing two; neuro/mental health is
the only empty band.}\end{minipage}

\vspace{2mm}
\noindent\begin{minipage}{\linewidth}\includegraphics[width=\linewidth]{\FIGDIR/f2_design_endpoint.png}
\figcap{Left --- architecture: four cohorts, four modelling/inventory records, one acute
time-series, one cross-sectional, one toxicology, one laboratory, one review; \textbf{metadata-only 5 (28\,\%)},
day-0 Elsevier and one ACS record. Right --- endpoints span AF, AMI, lipids, spirometry, frailty, cervical cancer
prognosis, HFMD and birthweight.}\end{minipage}

\vspace{2mm}
\noindent\begin{minipage}{\linewidth}\includegraphics[width=\linewidth]{\FIGDIR/f3_metric_geography.png}
\figcap{Left --- particle metric: PM$_{2.5}$-only 6, joint PM$_{2.5}$+PM$_{10}$ 4,
unspeciated/emissions 5, one composition record (sulfate), one optical (dust optical depth), one respirable dust.
Right --- geography: China 9 dominates; MENA 2; one each for the US, Southeast Asia, South Asia, Nigeria.}\end{minipage}

\vspace{2mm}
\noindent\begin{minipage}{\linewidth}\includegraphics[width=\linewidth]{\FIGDIR/f6_lifecourse.png}
\figcap{Life-course windows: working age leads (spirometry, lipids, cervical cancer, silica, AMI);
three older-adult records (Medicare AF, CHARLS frailty, AMI); one in utero (folic acid, metadata only), one childhood
(HFMD); no adolescent record.}\end{minipage}

\vspace{2mm}
\noindent\begin{minipage}{\linewidth}\includegraphics[width=\linewidth]{\FIGDIR/f4_heatmap.png}
\figcap{Subtopic $\times$ study architecture. Cardiovascular sits in cohorts and health-impact
modelling; exposure assessment and sensing fall mostly in metadata-only cells today, which limits what the
measurement side of this issue can assert.}\end{minipage}

\vspace{2mm}

\noindent\begin{minipage}{\linewidth}\includegraphics[width=\linewidth]{\FIGDIR/f5_forest.png}
\figcap{Ratio scale, four estimates from two studies: source-specific PM$_{2.5}$ (secondary sulfate, vehicle) and
incident AF per IQR in Medicare; PM$_{2.5}$ per IQR and cervical cancer recurrence (Fine--Gray SHR) and death in a
single-centre cohort. Per-IQR contrasts differ by source and study. Not poolable.}\end{minipage}

\vspace{3mm}

\noindent\begin{minipage}{\linewidth}\includegraphics[width=\linewidth]{\FIGDIR/f5_forest_2.png}
\figcap{Difference scale: FEV$_1$ change (mL) per IQR long-term PM$_1$, China Pulmonary Health Study
(n\,=\,50,991). Single estimate shown for scale; the coarse-fraction small-airway estimates are not given with
CIs in the abstract.}\end{minipage}

\vspace{4mm}

\par\vskip 0pt plus 24\baselineskip\penalty-200\vskip 0pt plus -24\baselineskip
\section*{\sffamily\bfseries\color{Deep}Paper-level digest}
\vspace{-1.5mm}{\color{Amber}\rule{\linewidth}{1.1pt}}\vspace{2.5mm}

{\fontsize{7.4}{9}\selectfont\color{Slate}Ordered by cluster. Each entry gives the
methodological core and the finding that matters, not an abstract paraphrase. Records
marked \textit{metadata only} carry no recoverable abstract and assert no finding.\par}

\band{Cardiovascular \& metabolic\hfill 3 records}{Teal}

\paperentry{Zhu source-specific PM$_{2.5}$ atrial fibrillation Medicare cohort 2026}
{Source-Specific PM$_{2.5}$ Exposure and Atrial Fibrillation Risk in U.S. Older Adults: A Nationwide Cohort Study}
{Environ Health (ACS)}
{26.3 million Medicare beneficiaries, 2000-2019, 171.6 M person-years, 4.44 M incident AF. Per IQR of source-specific PM$_{2.5}$: secondary sulfate HR 1.056 (1.053-1.058), heavy oil/coal combustion 1.048 (1.046-1.050), secondary nitrate 1.032 (1.030-1.034), vehicle 1.029 (1.028-1.030); biomass burning and dust modest, SOA null overall but positive in the West. Limitation: source profiles come from sparse speciation monitors interpolated by inverse distance over 100 km, so source contrasts are largely regional contrasts; sulfate and nitrate fields are smooth and the per-IQR scaling differs by source, so 'sulfate is most toxic' is not a per-microgram statement. Claims-based AF under-ascertains paroxysmal AF. Relevance: high - speciation, not mass, is what optical LCS cannot give; source-resolved risk raises the value of co-located filter or ACSM nodes.}
{10.1021/envhealth.5c00638}

\paperentry{Sun long-term air pollution lipid profiles Shanxi repeated measures 2026}
{Long-term exposure to ambient air pollution and lipid profiles: a longitudinal study in a heavily industrialized region of northern China}
{Front Public Health}
{Repeated-measures study of 9,040 adults with 19,113 lipid measurements from a health-screening programme, 2019-2023. Per 10 $\mu$g\,m$^{-3}$ PM$_{2.5}$: TG +1.85\%, TC +1.70\%, LDL-C +1.81\%, HDL-C -1.55\% (linear mixed models); PM$_{10}$, NO$_2$, SO$_2$ and O$_3$ in the same direction, CO weaker; age modified the HDL-C association after FDR correction. Limitation: a health-screening population self-selects; CHAP is itself trained on the sparse Shanxi monitor network that underrepresents coal-industry hotspots, and with five correlated pollutants the 'PM$_{2.5}$ most pronounced' ranking is not identified from single-pollutant models. Lipids are sensitive to fasting status, not mentioned. Relevance: moderate - industrial-region exposure surfaces are where LCS densification would most change CHAP errors.}
{10.3389/fpubh.2026.1950282}

\paperentry{Du Air Health Index acute myocardial infarction 325 cities 2026}
{Establishing an Air Health Index for Predicting the Risk of Acute Myocardial Infarction Onset Associated with Air Pollution and Ambient Temperature}
{Environ Sci Technol}
{1,207,625 AMI cases across 325 Chinese cities; excess risks from PM$_{2.5}$, NO$_2$, O$_3$ and temperature aggregated into a normalised AHI. Per IQR, AHI associated with +5.97\% (5.32-6.63\%) AMI incidence vs +3.33\% (2.93-3.72\%) for AQHI and +2.04\% (1.73-2.35\%) for AQI; near-linear, no threshold; validated on an independent external dataset; 8,534 per 100,000 AMI cases attributed. Limitation: the index is fitted to the same outcome it is scored on, so in-sample superiority over AQI/AQHI is expected; IQR-scaled comparisons across indices with different distributions are not like-for-like, and 'attributable' fractions for a composite that includes temperature mix avoidable and non-avoidable risk. Relevance: high - an outcome-weighted index is a natural output layer for real-time LCS networks, provided the PM$_{2.5}$ input is calibrated.}
{10.1021/acs.est.6c07635}

\par\vskip 0pt plus 10\baselineskip\penalty-200\vskip 0pt plus -10\baselineskip
\band{Respiratory \& allergic\hfill 1 record}{Sky}

\paperentry{Li size-specific PM lung function China Pulmonary Health Study 2026}
{Associations between Size-Specific Particulate Matter and Lung Function: Results from the China Pulmonary Health Study}
{Environ Health (ACS)}
{50,991 adults >=20 y from the China Pulmonary Health Study; multivariable linear models on long-term size-resolved PM. Per IQR PM$_{1}$: FEV1 -36.84 mL (-45.91 to -27.76), FVC -34.63 mL (-45.48 to -23.79); large-airway decrements grew as particle size fell, whereas the coarse fraction PM$_{2.5}$-10 gave the largest small-airway (FEF75, FEF25-75) decrements. Limitation: size fractions are derived from correlated modelled surfaces (PM$_{1}$ and PM$_{2.5}$ share most of their mass), and the abstract reports no multi-fraction model, so 'size-specific' ranking may reflect differential exposure-model error rather than biology; cross-sectional. Relevance: high - optical LCS report PM$_{1}$ and coarse bins natively, and this is evidence that those channels carry distinct health signal worth calibrating, not discarding.}
{10.1021/envhealth.5c00676}

\par\vskip 0pt plus 10\baselineskip\penalty-200\vskip 0pt plus -10\baselineskip
\band{Reproductive \& developmental\hfill 1 record}{Sage}

\paperentry{Ge preconception folic acid air pollution birthweight 2026}
{Preconception Folic Acid Supplementation Associated with Attenuated Effects of Air Pollution on Birthweight Outcomes}
{Environ Health (ACS)}
{\textsc{metadata only.} Title and authors only (Ge, Zhang, Liao ... Berhane, Bai et al.); no abstract retrievable from PubMed, Europe PMC, Crossref or Semantic Scholar at harvest. The title claims effect modification of the air pollution-birthweight association by preconception folic acid use. Pollutants, cohort, n and effect sizes cannot be stated. Limitation: unverifiable beyond the title; supplement use is self-reported and strongly socially patterned, so 'attenuation' by folate is exposed to confounding by SES and antenatal-care quality unless handled explicitly. Relevance: low - no sensing content evident.}
{10.1021/envhealth.5c00779}

\par\vskip 0pt plus 10\baselineskip\penalty-200\vskip 0pt plus -10\baselineskip
\band{Mechanistic toxicology\hfill 1 record}{Violet}

\paperentry{Liu PM$_{2.5}$ IL-17 microglia neuroinflammation 2026}
{IL-17 Signaling Pathway: A Potential Signaling Pathway Mediating Neuroinflammation Induced by PM$_{2.5}$}
{Environ Health (ACS)}
{BV2 microglia exposed to PM$_{2.5}$ lost viability and raised inflammatory cytokines; microarray DEGs enriched for inflammatory GO terms and the KEGG IL-17 pathway with downstream NF-kB and MAPK. An anti-IL-17A antibody reversed cytokine release, microglial activation and behavioural deficits (implying an in vivo arm). Limitation: the abstract gives no PM$_{2.5}$ source, composition or dose, so environmental relevance cannot be judged; BV2 is an immortalised murine line with an exaggerated inflammatory phenotype, and IL-17 is usually T-cell/gamma-delta derived, so microglial IL-17 signalling needs cell-type attribution. Relevance: low - mechanistic only.}
{10.1021/envhealth.6c00080}

\par\vskip 0pt plus 10\baselineskip\penalty-200\vskip 0pt plus -10\baselineskip
\band{Exposure assessment \& modelling\hfill 3 records}{Clay}

\paperentry{Mai population-weighted PM$_{2.5}$ O$_3$ China 2015-2024 preprint 2026}
{Population Weighted Ambient PM$_{2.5}$ and O$_3$ Concentrations in China from 2015 to 2024}
{Research Square (preprint)}
{Ecological geospatial analysis of annual \textasciitilde{}1 km CHAP PM$_{2.5}$/O$_3$ rasters matched to WorldPop counts, 2015-2024. Population-weighted PM$_{2.5}$ fell from 53.06 to 31.25 $\mu$g\,m$^{-3}$ (-41.1\%) while population-weighted O$_3$ rose from 77.06 to 96.75 $\mu$g\,m$^{-3}$ (+25.6\%); the population-weighted minus geographic gap narrowed for PM$_{2.5}$ (5.77 to 2.43) and widened for O$_3$ (1.01 to 2.72). 97.4\% of 0.2-degree blocks show falling PM$_{2.5}$ with rising O$_3$ at a 0.5 $\mu$g\,m$^{-3}$/y threshold. Limitation: purely descriptive and inherits CHAP's ML smoothing; no uncertainty on the weighted means; preprint, not peer reviewed. Relevance: moderate - the narrowing urban-rural PM$_{2.5}$ gap shifts LCS value toward peri-urban and rural siting.}
{10.21203/rs.3.rs-10964316/v1}

\paperentry{Tripathi satellite-reanalysis sulfate aerosol India ground measurements 2026}
{Improving Satellite-Reanalysis Sulfate Aerosol Estimates over India Using Ground Measurements}
{Environ Pollut}
{\textsc{metadata only.} Title and authors only (Tripathi, Sahu, Sharma, Gupta); no abstract indexed at harvest. The title indicates ground measurements used to correct satellite-reanalysis (likely MERRA-2/CAMS-type) sulfate aerosol over India. Data sources, method and skill cannot be stated. Limitation: unverifiable beyond the title; Indian speciation networks are sparse and urban-biased, so any correction trained on them risks transferring urban bias to rural cells. Relevance: moderate - component-resolved surfaces are the complement LCS mass networks lack; pairs with the sulfate-dominant AF signal in today's Medicare cohort.}
{10.1016/j.envpol.2026.129391}

\paperentry{Dayanandana CMIP6 dust optical depth MENA evaluation 2026}
{Evaluation of Dust Optical Depth in CMIP6 Simulations over the Middle East-North Africa Region}
{Atmos Environ}
{\textsc{metadata only.} Title and authors only (Dayanandana, Parottil, Abbas, Al-Harrasi et al.); no abstract at harvest. Evaluates CMIP6 dust optical depth over MENA, presumably against MODIS/MERRA-2 DOD. Models, metrics and biases cannot be stated. Limitation: unverifiable beyond the title; column optical depth maps to surface PM$_{10}$ only through assumed vertical profiles and size distributions, so climate-model DOD skill says little about surface exposure. Relevance: low-moderate - MENA dust is where optical LCS mass error from coarse particles is largest.}
{10.1016/j.atmosenv.2026.122431}

\par\vskip 0pt plus 10\baselineskip\penalty-200\vskip 0pt plus -10\baselineskip
\band{Sensing, forecasting \& instrumentation\hfill 2 records}{Amber}

\paperentry{Pyae Tin Saw portable PM$_{2.5}$ cross-site prediction fixed station 2026}
{Short-term PM$_{2.5}$ data continuity: Cross-site prediction of portable measurements using a nearby fixed monitoring station}
{Environ Pollut}
{\textsc{metadata only.} Title and authors only (Pyae Tin Saw, Kallawicha); day-0 Elsevier deposit with no abstract anywhere at harvest. The title describes predicting short-term portable PM$_{2.5}$ records from a nearby fixed monitoring station to maintain data continuity. Sensor model, distance, n, and skill metrics cannot be stated. Limitation: unverifiable beyond the title; cross-site transfer functions are typically valid only within the regional-background regime and fail for local sources, which is exactly when portable data matter. Relevance: high - gap-filling an LCS or portable record against a regulatory anchor is a core network-operations problem.}
{10.1016/j.envpol.2026.129384}

\paperentry{Ghaffarpasand Pope vehicle remote sensing real-world emissions high emitters 2026}
{Emerging Evidence on Real-World Exhaust Emissions and High Emitting Vehicles from Remote Sensing Observations}
{Atmos Pollut Res}
{\textsc{metadata only.} Title and authors only (Ghaffarpasand, Pope); no abstract at harvest. Appears to synthesise vehicle emission remote-sensing (VERSS) evidence on real-world exhaust and high emitters. Whether PM/smoke opacity or only NOx/CO/HC is covered cannot be stated. Limitation: unverifiable beyond the title; VERSS PM is inferred from UV opacity with large per-pass uncertainty, and single-pass snapshots overweight hot-running conditions. Relevance: moderate - high-emitter fractions explain the kerbside spikes roadside LCS record but cannot attribute.}
{10.1016/j.apr.2026.103244}

\par\vskip 0pt plus 10\baselineskip\penalty-200\vskip 0pt plus -10\baselineskip
\band{Occupational \& indoor\hfill 1 record}{Amber}

\paperentry{Alquwayi respirable crystalline silica occupational exposure Saudi Arabia review 2026}
{Potential risk of occupational exposure to respirable crystalline silica in Saudi Arabia: a comprehensive review of the current status and future challenges}
{Front Public Health}
{Single-author review integrating literature, labour and mining statistics, mineralogy and regulation through a 'silica hazard triangle' (workforce, materials, activities). Concludes that mining, construction, cement, glass, engineered stone and photovoltaic manufacturing expansion plus arid-climate dust raise RCS exposure potential, while measured exposure, epidemiology and surveillance data are scarce. Limitation: no new measurements and no quantitative exposure estimates - the conclusion is an inference from sector size, not from data. Relevance: low-moderate - optical LCS cannot speciate silica, so RCS monitoring needs gravimetric/FTIR sampling regardless.}
{10.3389/fpubh.2026.1918362}

\par\vskip 0pt plus 10\baselineskip\penalty-200\vskip 0pt plus -10\baselineskip
\band{Burden, policy \& mitigation\hfill 3 records}{Coral}

\paperentry{Zhang respirable magnetite emissions coal-fired power China 2026}
{Emissions of Respirable Magnetite from the Coal-Fired Power Sector in China: Spatial Inequalities and Mitigation Pathways}
{Environ Sci Technol}
{Multifactor emission model for respirable magnetite from Chinese coal-fired power: 2,828.7 (2,307.1-3,371.4) t/y, spatially concentrated by tiered emission standards. Guizhou and Sichuan supply 4\% of coal power but tightening their standards alone would cut >31\% of national magnetite emissions. Limitation: emission factors for magnetite specifically (as opposed to total PM or Fe) rest on few stack measurements, and the uncertainty interval looks narrow for a species-level inventory; health relevance is asserted from brain-magnetite literature rather than exposure-response data. Relevance: moderate - magnetite is invisible to optical LCS mass; magnetic or filter-based speciation would be needed to verify hotspots.}
{10.1021/acs.est.6c08602}

\paperentry{Idowu medical waste incinerator airborne heavy metals Monte Carlo Nigeria 2026}
{Monte Carlo-Based Probabilistic Health Risk Assessment of Airborne Heavy Metals from a Medical Waste Incinerator in Southwestern Nigeria}
{Research Square (preprint)}
{Particulate matter collected on passive glass slides around a medical-waste incinerator for two months; ten metals analysed, converted to airborne concentrations, and propagated through 10,000 Monte Carlo draws. Median HQ \textasciitilde{}400, 95th percentile 850; exceedance probability near 100\%; ILCR above benchmarks. Limitation: passive deposition slides measure settled mass, not airborne concentration - the conversion to C\_air requires an assumed deposition velocity that dominates the result, and an HQ of 400 almost certainly reflects that conversion rather than real inhalation dose. Preprint. Relevance: moderate as a cautionary example - fence-line LCS or low-volume samplers would have given a defensible air concentration.}
{10.21203/rs.3.rs-11381994/v1}

\paperentry{Jiang PLA melt-blown nonwoven filtration durability humidity airflow 2026}
{Service Durability of PLA Melt-Blown Nonwovens: Roles of Humidity and Airflow Velocity in Degradation and Filtration Performance}
{Langmuir}
{Full-factorial service-condition study (humidity, airflow velocity, PM loading, natural microbes) on PLA melt-blown filter media. Degradation stays early-stage, first-order, k up to 3.88e-2 /h, surface-erosion dominated. Humidity drives degradation (90\% of short-term water-contact-angle variance; 77\% of yield-strength and 91\% of pressure-drop variance long-term); airflow velocity governs dynamic filtration via particle-media contact. Limitation: the abstract gives no filtration efficiency numbers by particle size, so it is unclear whether MPPS performance falls within a realistic mask or HVAC service life. Relevance: low - mitigation materials; humid-climate filter life matters for purifier intervention trials.}
{10.1021/acs.langmuir.6c03984}

\par\vskip 0pt plus 10\baselineskip\penalty-200\vskip 0pt plus -10\baselineskip
\band{Other clinical endpoints\hfill 3 records}{Slate}

\paperentry{Lin air pollution cervical cancer recurrence mortality Fudan cohort 2026}
{Long-Term Exposure to Ambient Air Pollution and Cervical Cancer Outcomes in Chinese Women: A Prospective Study on Recurrence and Mortality Risks}
{Environ Health (ACS)}
{3,758 women with primary cervical cancer (2019-2022), median follow-up \textasciitilde{}3 y; 217 recurrences, 113 deaths. Per IQR PM$_{2.5}$: recurrence SHR 1.93 (1.40-2.66, Fine-Gray), death HR 2.55 (1.78-3.67); PM$_{10}$ SHR 1.65 (1.13-2.42) and HR 2.21 (1.38-3.54); NO$_2$ and SO$_2$ also positive, O$_3$ null. Limitation: hazard ratios near 2-2.5 per IQR from 113 deaths are implausibly large for a residential long-term contrast and point to confounding by referral geography - patients from outside Shanghai (higher PM) likely differ in stage, HPV management and follow-up intensity; single-centre, short follow-up, no adjustment for distance to the hospital reported in the abstract. Relevance: low - model-based residential exposure, no sensing implication.}
{10.1021/envhealth.6c00011}

\paperentry{Lin PM$_{2.5}$ frailty trajectories CHARLS 2026}
{The association of particulate matter 2.5 exposure with frailty trajectories in middle-aged and older Chinese adults and the potential modification by lifestyle and socioeconomic status}
{J Occup Environ Med}
{10,659 CHARLS participants; deficit-accumulation frailty index, three GBTM trajectories (stable low 64\%, increasing to frailty 29\%, continuously worsening 7\%). Per 10 $\mu$g\,m$^{-3}$ PM$_{2.5}$, odds of the two adverse trajectories rose 15.2\% and 12.3\% (ordered logistic; no CIs in the abstract). Higher SES attenuated the association (RERI -0.067); unhealthy sleep duration amplified it (RERI 0.163). Limitation: GBTM classes are model artefacts with posterior misclassification carried into a second-stage regression without uncertainty propagation; CHARLS addresses are coarse (city/county), so exposure contrast is mostly between regions with different healthcare and diet. Relevance: low - regional satellite exposure, no sensing implication.}
{10.1097/jom.0000000000003913}

\paperentry{Jiang air pollution meteorology hand foot mouth disease Shenzhen DLNM 2026}
{Ambient air pollution and meteorological factors associated with hand, foot, and mouth disease risk in Shenzhen, China: a time-series study}
{Front Public Health}
{Shenzhen daily HFMD counts 2013-2022 with six pollutants and meteorology in DLNMs, including pre/post EV71 vaccine and COVID-19 periods. PM$_{2.5}$ and PM$_{10}$ showed risk peaks at low concentrations and 'protective' effects at extreme highs; SO$_2$ and CO raised cumulative risk at extremes; temperature showed a dual-peak pattern and high humidity raised risk. No effect sizes with CIs in the abstract. Limitation: HFMD is a contact-transmitted enteroviral disease with strong seasonality driven by school terms; non-monotone, sign-flipping PM curves are the signature of residual seasonal confounding and DLNM over-flexibility, not biology. Relevance: low.}
{10.3389/fpubh.2026.1920424}

\clearpage
\section*{\sffamily\bfseries\color{Deep}Corpus register}
\vspace{-1.5mm}{\color{Amber}\rule{\linewidth}{1.1pt}}\vspace{2.5mm}

\input{table_rows}

\vspace{2mm}

\begin{keybox}
{\sffamily\bfseries\footnotesize\color{Deep}Provenance and method}\\[1.2mm]
{\fontsize{7.2}{8.8}\selectfont
Entry window \textbf{21 October 2026}, single day (Wednesday). Harvester legs: \textbf{PubMed} health
\textbf{2}, sensing \textbf{0} (connector EDAT query confirms 2 total --- indexing lull); \textbf{Europe PMC}
\textbf{53} (2 AGRICOLA dropped); \textbf{Crossref by ISSN} \textbf{61}; \textbf{arXiv} HTTP 429;
\textbf{OpenAlex} HTTP 429. 113 unique, 10 already carried, \textbf{103 fresh}; PubMed connector (EDAT 21 Oct)
2 PMIDs, 0 appended. Europe PMC records lacking DOI were resolved by Crossref title match (five returned
supplementary-file \texttt{.s001} DOIs, stripped and verified) and their abstracts recovered from PubMed by DOI.
\textbf{Consensus} calibration sweep: 10 hits, 1 candidate (Ginsburg 2026, PurpleAir global calibration) ---
\textbf{rejected by the DOI gate as already shipped 29 July}. \textbf{104} candidates screened, \textbf{18
admitted}, \textbf{86 rejected} with logged reasons. \textbf{DOI gate} 0 fail / 0 warn after the duplicate was
removed. \textbf{Trial watch} 21--22 October: 2 registrations updated (AIRWEIGHS purifier RCT NCT02763917, primary
completion passed; Taif COPD survey). \textbf{5 carried on metadata alone}. \textbf{Scholar Gateway} not called.
Abstracts remain the copyright of their publishers.\par}
\end{keybox}

\end{document}
